\chapter{导数、微分与高阶展开}
\section{基本公式与求导结构}
导数是线性近似的系数：$f(a+h)=f(a)+f'(a)h+o(h)$。可导必连续，连续未必可导，如 $|x|$ 在 0。左右导数相等且有限才可导。微分 $\dd y=f'(x)\dd x$；复合函数的一阶微分形式不变，高阶微分一般不能直接照搬。
\begin{longtable}{p{.44\textwidth}p{.45\textwidth}}
\toprule 函数与导数 & 条件或补充\\\midrule\endhead
$(x^\alpha)'=\alpha x^{\alpha-1}$ & 一般实指数取 $x>0$，整数次幂可按其定义域扩展\\
$(\ee^x)'=\ee^x$；$(a^x)'=a^x\ln a$ & $a>0$\\
$(\ln|x|)'=1/x$；$(\log_a x)'=1/(x\ln a)$ & 前者 $x\ne0$；后者 $x>0,a>0,a\ne1$\\
$(\sin x)'=\cos x$；$(\cos x)'=-\sin x$ & $x\in\R$\\
$(\tan x)'=\sec^2x$；$(\cot x)'=-\csc^2x$ & 避开原函数无定义处\\
$(\sec x)'=\sec x\tan x$ & $(\csc x)'=-\csc x\cot x$\\
$(\arcsin x)'=1/\sqrt{1-x^2}$ & $(\arccos x)'=-1/\sqrt{1-x^2}$，$|x|<1$\\
$(\arctan x)'=1/(1+x^2)$ & $(\operatorname{arccot}x)'=-1/(1+x^2)$，采用常见主值约定\\
$(\sinh x)'=\cosh x$ & $(\cosh x)'=\sinh x$；$(\tanh x)'=1/\cosh^2x$\\
$(\arsinh x)'=1/\sqrt{1+x^2}$ & $\arsinh x=\ln(x+\sqrt{1+x^2})$\\\bottomrule
\end{longtable}
\ruleitem{线性、乘积、商与复合}{
$(uv)'=u'v+uv'$；$(u/v)'=(u'v-uv')/v^2$（$v\ne0$）；$(f\circ g)'=f'(g)g'$。
乘积与复合交错出现时，先标清最外层运算。}
\ruleitem{对数求导}{若 $y=u(x)^{v(x)}$ 且 $u>0$，则 $y'=u^v[v'\ln u+vu'/u]$。多因子乘除也可先求 $\ln|y|$，再还原 $y'$。}
\ruleitem{反函数}{若 $y=f(x)$ 在某区间可逆且 $f'(x)\ne0$，则 $(f^{-1})'(y)=1/f'(x)$，二阶为 $-f''(x)/(f'(x))^3$。每次求导要说明自变量是谁。}
\ruleitem{参数方程}{若 $x=x(t),y=y(t)$ 且 $x'(t)\ne0$，
\[
\frac{\dd y}{\dd x}=\frac{y'}{x'},\qquad
\frac{\dd^2y}{\dd x^2}=\frac{x'y''-y'x''}{(x')^3}.
\]}
\ruleitem{一元隐函数}{对 $F(x,y)=0$，若 $F\in C^1$ 且 $F_y\ne0$，局部存在 $y=y(x)$，
\[
y'=-\frac{F_x}{F_y},\qquad
y''=-\frac{F_{xx}+2F_{xy}y'+F_{yy}(y')^2}{F_y}\quad(F\in C^2).
\]}
\ruleitem{高阶乘积公式}{若相关导数存在，
$\displaystyle(uv)^{(n)}=\sum_{k=0}^n\binom nk u^{(k)}v^{(n-k)}$。
常用 $\frac{\dd^n}{\dd x^n}\ee^{ax}=a^n\ee^{ax}$，$\frac{\dd^n}{\dd x^n}\sin(ax+b)=a^n\sin(ax+b+n\pi/2)$。}
\ruleitem{系数提取}{若 $f(a+h)=\sum_{k=0}^n c_kh^k+o(h^n)$ 且相应可导，则 $f^{(n)}(a)=n!c_n$。不是直接取 $c_n$。}

\section{泰勒公式与几何应用}
\ruleitem{Peano 余项}{在 $a$ 具有 $n$ 阶导数时，
$\displaystyle f(a+h)=\sum_{k=0}^n\frac{f^{(k)}(a)}{k!}h^k+o(h^n)$。
这是局部极限工具，不能据此直接估算远处误差。}
\ruleitem{Lagrange 余项}{若 $f$ 在 $a,x$ 之间有到 $n+1$ 阶的充分光滑性（例如 $C^{n+1}$），则
\[
f(x)=\sum_{k=0}^n\frac{f^{(k)}(a)}{k!}(x-a)^k+
\frac{f^{(n+1)}(\xi)}{(n+1)!}(x-a)^{n+1},
\]
$\xi$ 在 $a,x$ 之间；若 $|f^{(n+1)}|\le M$，误差绝对值不超过 $M|x-a|^{n+1}/(n+1)!$。}
\ruleitem{积分余项}{若 $f\in C^{n+1}$，余项也可写为
$\displaystyle R_n(x)=\frac1{n!}\int_a^x f^{(n+1)}(t)(x-t)^n\dd t$。
带权积分不等式常比不定的 $\xi$ 更方便。}
\ruleitem{切线、法线}{曲线 $y=f(x)$ 在 $a$ 的切线 $y-f(a)=f'(a)(x-a)$；法线用 $(x-a)+f'(a)(y-f(a))=0$，此式也适用于 $f'(a)=0$。参数曲线可用方向向量 $(x',y')$ 写切线，避免遗漏竖直切线。}
\ruleitem{弧长与曲率}{
\[
\dd s=\sqrt{1+(y')^2}\dd x,quad
\kappa=\frac{|y''|}{(1+(y')^2)^{3/2}},\quad
\kappa=\frac{|x'y''-y'x''|}{((x')^2+(y')^2)^{3/2}}.
\]
后式要求曲线正则，即 $(x',y')\ne(0,0)$。曲率半径为 $1/\kappa$（$\kappa\ne0$）。}
\ruleitem{数值求根（补充）}{牛顿迭代 $x_{n+1}=x_n-f(x_n)/f'(x_n)$。若 $f\in C^2$、$f(r)=0$、$f'(r)\ne0$，初值足够靠近简单根 $r$，则局部二阶收敛；远离根时可能发散，不能无条件使用。}

\section{高阶导数与隐函数真题}
\problem{15}{指数与正弦乘积的四阶导数}{2016 年第 8 届，非数学类，一(4)，\src{08}}
求 $f(x)=\ee^x\sin2x$ 的 $f^{(4)}(0)$。
\method{只取 $x^4$ 的系数，比连续四次求导省算。}
\sol $\ee^x=1+x+x^2/2+x^3/6+O(x^4)$，$\sin2x=2x-4x^3/3+O(x^5)$。
乘积的 $x^4$ 系数为 $2/6-4/3=-1$，故 $\ans{f^{(4)}(0)=-24}$。
也可用复数 $(1+2\ii)^4=-7-24\ii$ 复核虚部。部分网络参考解把符号写为正，这里以计算结果为准。

\problem{16}{零点重数与莱布尼茨公式}{2020 年第 12 届，非数学类，一(2)，\src{12}}
求 $f(x)=(x+1)^n\ee^{-x^2}$ 的 $f^{(n)}(-1)$（$n\ge1$）。
\sol 对 $(x+1)^n$ 求不足 $n$ 次导数，代入 $-1$ 均为 0。莱布尼茨和中只剩该因子求 $n$ 次、指数因子不求导的一项，故 $\ans{n!/\ee}$。等价地令 $h=x+1$，$f(-1+h)=\ee^{-1}h^n+O(h^{n+1})$。

\problem{17}{部分分式给出任意阶导数}{2023 年第 15 届，非数学 A 类，一(3)，\src{15A}}
设 $f(x)=1/(x^2-3x+2)$，求 $f^{(n)}(0)$。
\sol $f(x)=1/(1-x)-1/(2-x)$，在 $|x|<1$ 展开得
\[f(x)=\sum_{n=0}^\infty\left(1-\frac1{2^{n+1}}\right)x^n.\]
故 $\ans{f^{(n)}(0)=n!(1-2^{-n-1})}$。先分解两个简单极点，避免直接对有理函数反复求导。

\problem{18}{含变上限积分的二阶隐导数}{2024 年第 16 届，非数学 B 类，一(3)，\src{16B}}
$\ee^{3y}+\int_0^{x+y}\cos(t^2)\dd t=1$ 确定过原点的隐函数，求 $y''(0)$。
\sol 一次求导：
\[3\ee^{3y}y'+\cos((x+y)^2)(1+y')=0.\]
代入原点得 $4y'(0)+1=0$。再求导时，$\cos((x+y)^2)$ 的导数为
$-2(x+y)\sin((x+y)^2)(1+y')$，原点值为 0，所以
\[9(y'(0))^2+4y''(0)=0,\qquad \ans{y''(0)=-9/64}.\]
切勿漏掉 $\ee^{3y}$ 求导得到的 $(y')^2$ 项。

\problem{19}{参数曲线与隐函数合用}{2018 年第 10 届，非数学类，一(2)，\src{10}}
$x=t+\cos t$，$\ee^y+ty+\sin t=1$，求 $t=0$ 时的切线。
\sol $t=0$ 时 $(x,y)=(1,0)$，且 $x'(0)=1$。第二式对 $t$ 求导：
$(\ee^y+t)y'+y+\cos t=0$，故 $y'(0)=-1$。切线为 $\ans{x+y-1=0}$。

\problem{20}{小量恒等式同时确定函数值与导数}{第 14 届初赛补赛二，一(2)，\src{14}}
在 1 附近可导的 $f$ 满足
$f(1+x)-3f(1-x)=4+2x+o(x)$，求曲线在 $x=1$ 处切线。
\sol 左端展开为 $-2f(1)+4f'(1)x+o(x)$。比较常数、一次项得 $f(1)=-2$、$f'(1)=1/2$，所以 $y+2=(x-1)/2$，即 $\ans{x-2y-5=0}$。
“比较系数”在此可严格化：先令 $x\to0$，再减去常数项、除以 $x$ 求极限。

\problem{21}{零点重数让复合求导变简单}{2025 年第 17 届，非数学 B 类，一(1)，\src{17B}}
$y=x^2\ee^{\sin x}$，求 $y''(0)$。
\sol $\ee^{\sin x}=1+O(x)$，所以 $y=x^2+O(x^3)$；函数光滑，由二阶系数得 $\ans{y''(0)=2}$。若用莱布尼茨公式，也仅有 $(x^2)''\ee^{\sin0}$ 一项保留。
